%% Introduction %%

Cela dit, le gros bémol que cet outil présente est la difficulté à le situer dans le temps. Pour le moment, cet outil figure parmi les projets du département de la direction des patrimoines. Son développement ne prendra place qu'à la direction Data \& Technologies dans une temporalité que nous ignorons. Il devient alors difficile d'évaluer quand débutera son développement et en conséquent de quels outils actuellement en développement, réflexion et en tests les utilisateurs des collections de l'INA disposeront-ils ? Les besoins peuvent varier si par exemple les chercheurs et professionnels de l'audiovisuel disposent des transcriptions des émissions ou encore de si les résumés par IA ont pleinement intégrés les bases de données. 

\subsection{Une question de base de données}

Le point sur lequel l'outil peut particulièrement varier est la base de données dans laquelle cet outil ira récupérer ses données. En l'état actuel, il a été pensé et construit sur l'exploitation des bases de données du dépôt légal. Si le lac de données vient à être en place pour plusieurs des publics visés, les défauts de qualités des données peuvent avoir été corrigés par cette nouvelle base de données. Il sera nécessaire d'adapter les datavisualisation pour le nouveau modèle de données. Les informations récupérées ne seront plus exactement les mêmes. Nous avions par exemple imaginés faire appel à wikidata afin de faire apparaître quelle chaîne succédait à quelle chaîne, sous quel nom. Ce que la base de données du dépôt légal ne permettait pas de faire mais que le nouveau modèle de données permettrait. 

D'un autre côté, il est pour le moment difficile de se projeter dans le lac de données car pour le moment, les données du dépôt légal ne sont pas intégrés à cet outil, mais seulement ses référentiels. 

\subsection{Manque de pertinence pour tous les publics}

Cet outil manque de pertinence pour les professionnels de l’audiovisuel qui seront confronté avec cet outil à des documents sur lesquels l'INA n'a aucun droit d'exploitation sans que ces indications ne paraissent dans les fiches créées par cet outil. Par nature, cet outil est conçu pour comprendre et suivre le traitement documentaire du dépôt légal. Les recherches sont effectuée par le titre de l'émission, de sa collection ou de sa tranche horaire. Ses fonctionnalités de correspondent pas aux besoins de professionnels de l'audiovisuel. Ces derniers vont donc continuer à se reposer sur les documentalistes. 

%% Conclusion %%

C'est pourquoi cet outil a été pensé pour s'adapter à toutes les évolutions possibles. Les datavisualisations ont été conçu pour montrer certaines informations. Le nom que prend les champs affichés et le modèle de données que suit la base de données ne sont pas de la plus grand importance. Nous avons conçu un outil fait pour fournir un service. Il est conçu pour le faire et s'adapter, peut importe les changements. 